\documentclass[11pt,a4paper]{article}
\usepackage[utf8]{inputenc}
\usepackage[T1]{fontenc}
\usepackage{amsfonts}
\usepackage[french]{babel}
\frenchbsetup{StandardLists=true} % à inclure si on utilise \usepackage[francais]{babel}
\usepackage{enumitem}
\usepackage{amssymb}
\usepackage{amsmath}
\usepackage[left=2cm,right=2cm,top=2cm,bottom=2cm]{geometry}
\usepackage{multirow}
\usepackage[dvipsnames]{xcolor}

\usepackage[T1]{fontenc}
\usepackage{fouriernc}
\usepackage{graphicx}
\usepackage{colortbl}

\usepackage{setspace}
\singlespacing

\usepackage{pstricks}
\usepackage{xkeyval}
\usepackage{pst-labo}


\usepackage{multicol}
\usepackage{caption}
\usepackage{fancyhdr}
\pagestyle{fancy}
\renewcommand\headrulewidth{1pt}
\fancyhead[L]{Lycée Denis Diderot }
\fancyhead[R]{Classe de 1 Bac Pro }
\setcounter{page}{9}\fancyfoot[C]{page \thepage}
\fancyfoot[R]{Année 2024-2025}
\fancyfoot[L]{Alexis RUIZ}
\usepackage{multirow,tabularx,booktabs}
\newcommand{\cadreblanc}[1]{%
\medskip\framebox[\linewidth][c]{%
\rule{0mm}{#1mm}}\par
\medskip}

\usepackage{awesomebox}
\setlength{\aweboxvskip}{0mm}
\usepackage{pifont}
\usepackage{chemmacros}
\chemsetup{modules=all}
\setlength{\columnseprule}{.25mm}

\renewcommand{\arraystretch}{2}
\newcolumntype{C}[1]{>{\centering\arraybackslash }b{#1}}

\newcommand*\acocher{\quitvmode{\fboxsep0pt \fboxrule0.8pt \fbox{\vrule height1.8ex depth.1ex width0pt \vrule height0pt depth0pt width1.9ex }}}

\newcommand{\code}{\awesomebox[black]{2pt}{

\includegraphics[scale=0.15]{Chapitre 1 - Energie et puissance électrique/images/frame.png}
}{black}}

\newcommand{\moteur}{\awesomebox[black]{2pt}{

\includegraphics[scale=1]{images/moteur.jpg} 
}{black}}

\newcommand{\defconclusion}{\awesomebox[black]{2pt}{\rotatebox{90}{\Large{\textbf{Conclusion}}}}{black}}


  \newenvironment{itemize*}%  % listes compactes
    {\begin{itemize}%
      \setlength{\itemsep}{0pt}%
      \setlength{\parskip}{0pt}}%
    {\end{itemize}}


\frenchbsetup{StandardLists=true}
\newlist{fleche}{itemize}{1}
\setlist[fleche]{label=\ding{220},font=\color{orange}}


\usepackage[tikz]{bclogo} 
\newcommand{\ctext}[1]
{\begin{bclogo}[couleur = white,logo= , arrondi = 0.1,barre=none]
{}
\vspace{-8mm}
~~
\textcolor{red}{#1}
~~
\end{bclogo}}

%%%%%%%%%%%%%%%%%%%%%%%%%%%%%%%%%%%%%%%%%%%%%%%%%%%%%%%%%%%



\begin{document}


\begin{center}
 \LARGE{COURS } \\[3mm]
\end{center}



\normalsize

\hrule\
\\[0,5cm]


\fcolorbox{black}{white}{\textbf{\ding{202} Chute libre - voir TP2}}\\
\vspace{-3mm}
\begin{itemize*}
\item La chute libre est le mouvement dans le \textbf{vide} d'un solide soumis \textbf{uniquement à l'attraction terrestre. }
\vfill 


\ctext{
Remarque : lorsque les forces de résistance de l'air et de la poussée d'Archimède sont négligeables, on admet qu'un solide en tombant dans l'atmosphère est en chute libre. 
}
\vfill 
\begin{minipage}{.54\textwidth}%
 \item

 \ctext{ Un solide en chute libre, lâché sans vitesse initiale, est en \textbf{mouvement rectiligne} et son accélération est \textbf{constante}. \\
        \item - son déplacement est en ligne droite et \textbf{vertical}. \\
       \item - dans le cas de la chute libre, l'accélération $a$ (en $\mathrm{m/s^2}$ est égale à $g$ l'accélération de la pesanteur 
} 

 \item \ctext{La vitesse instantanée $v$ du solide augment proportionnellement au temps : \\
 $$v=g \times t = 9,8 \times t$$}

\end{minipage}%
\hfill
\begin{minipage}{.40\textwidth}%
\includegraphics[width=\textwidth]{Chapitre 2 - Accélération et vitesse dans le cas d’un mouvement rectiligne/images/2 - Cours1.png}

\end{minipage}%

\vfill 

\item 
\noindent%
\begin{minipage}{.54\textwidth}%
 \ctext{La chute libre est un mouvement uniformément accéléré : à chaque seconde de chute libre, la vitesse du solide augmente de 9,8 m/s. L'accélération est donc de 9,8 m/s par seconde donc 9,8 $m/s^2$
 \vspace{1cm}
 Remarque : La vitesse en chute libre est indépendante de la masse du solide. 
 \vspace{1cm} }

\end{minipage}%
\hfill
\begin{minipage}{.40\textwidth}%
\includegraphics[width=\textwidth]{Chapitre 2 - Accélération et vitesse dans le cas d’un mouvement rectiligne/images/2 - Cours2.png}
\end{minipage}%
\end{itemize*}
\vfill

\newpage 

\fcolorbox{black}{white}{\textbf{\ding{203} Mouvement rectiligne d'accélération constante}}\\
\vspace{-3mm}
\begin{itemize*}

\item Mouvement rectiligne uniformément varié [MRUV]

\noindent%
\begin{minipage}{.35\textwidth}%
 \ctext{Lorsqu'un solide est en mouvement rectiligne uniformément varié : \\
 \begin{itemize}
     \item son déplacement se fait en ligne droite
     \item son accélération $a$ en $m/s^2$ est constante
     \item si a>0 alors la vitesse augmente. 
     \item si a<0 alors la vitesse diminue
     \item la vitesse initiale $v_0$ est la vitesse de départ à l'instant $t=0 s$
     \vspace{1cm}
     \item La relation entre vitesse et accélération est : \\
     \Large{$$v=a \times t + v_0$$}
 \end{itemize}}

\end{minipage}%
\hfill
\begin{minipage}{.55\textwidth}%
\fbox{\includegraphics[width=\textwidth]{Chapitre 2 - Accélération et vitesse dans le cas d’un mouvement rectiligne/images/2 - Cours3.png}} \\[5mm]

\fbox{\includegraphics[width=\textwidth]{Chapitre 2 - Accélération et vitesse dans le cas d’un mouvement rectiligne/images/2 - Cours4.png}}
\end{minipage}%

\vfill 

\item Mouvement rectiligne uniforme[MRU] (Voir TP1) \\ 
\begin{fleche}
\item L'accélération est nulle : $a=0 m/s^2$
\item La vitesse est constante : $v=v_O=constante$
\end{fleche}

\item Ordre de grandeur de vitesses et accélérations : \\



    \begin{tabular}{|c||c|c|c|c|}
         \hline
         \rowcolor{orange} Véhicule &  \multicolumn{2}{|c|}{Vitesse} & Accélération & Décélération\\
         \hline 
        Train & ~~~~~~$150~ km/h$~~~~~~~~ & ~~~~~~~~ $42~ m/s$ ~~~~~~~~ & ~~~~~~~~ $2 ~m/s^2$ ~~~~~~~~ & ~~~~~~~~ $-5~~m/s^2$~~~~~~~~ \\
        \hline
         Ascenseur& $11~ km/h$ & $3~ m/s$ & $2 ~m/s^2$ & $-2~~m/s^2$\\
         \hline
         Automobile & $130~ km/h$ & $36~ m/s$ & $3 ~m/s^2$ & $-8~~m/s^2$\\
         \hline
         Avion de ligne & $800~ km/h$ & $222~ m/s$ & $4 ~m/s^2$ & $-3~~m/s^2$\\
         \hline
         Moto & $130~ km/h$ & $36~ m/s$ & $7 ~m/s^2$ & $-6~~m/s^2$\\
         \hline 
    \end{tabular}





\vfill 
\end{itemize*}











\end{document}